\begin{table*}[t]
\centering
\caption{Geographic characteristics and meteorological measurement heights of the WBEA atmospheric monitoring stations included in the study.}
\label{tab:stations}
\begin{tabular*}{\textwidth}{@{\extracolsep{\fill}}llrrll@{}}
\toprule
Station & Name & Latitude ($^\circ$N) & Longitude ($^\circ$W) & Wind height (m) & Temperature height (m) \\
\midrule
AMS 1 & Fort McKay - Bertha Ganter & 57.1894 & 111.6406 & 10 & 2, 10 \\
AMS 2 & Mildred Lake & 57.0500 & 111.5641 & 10 & 2 \\
AMS 5 & Mannix & 56.9680 & 111.4821 & 20, 45, 75, 90 & 2, 20, 45, 75, 90 \\
AMS 6 & Patricia McInnes & 56.7514 & 111.4767 & 10 & 2 \\
AMS 7 & Athabasca Valley & 56.7334 & 111.3905 & 10 & 2 \\
AMS 9 & Barge Landing & 57.1982 & 111.5995 & 10, 20 & 2 \\
AMS 11 & Lower Camp & 57.0268 & 111.5008 & 10 & 2 \\
AMS 13 & Fort McKay South & 57.1492 & 111.6423 & 10 & 2 \\
AMS 14 & Anzac & 56.4489 & 111.0380 & 20 & 2 \\
AMS 18 & Stony Mountain & 55.6215 & 111.1728 & 10, 20 & 2 \\
AMS 23 & Fort Hills & 57.3489 & 111.6397 & 10 & 2 \\
AMS 30 & Ells River & 57.2412 & 111.7221 & 10 & 2 \\
\bottomrule
\end{tabular*}
\par\smallskip\begin{minipage}{\textwidth}\footnotesize
Coordinates (WGS84) from WBEA 2024 site documents. Heights (m above ground) list variables with at least 5\% valid hourly values in the 2021--2025 exports; where 10 and 20~m are both listed, the 10~m sensor is used when valid. Methane, total hydrocarbons and non-methane hydrocarbons are reported together at every station (ppm, hourly) by WBEA's continuous hydrocarbon analysers, with an automated zero/span check that occupies one hour per day and is flagged in the data (WBEA SOP-ANA-008; WBEA monthly reports). The analyser model is documented here only for AMS~5 (Thermo~55i, installed 2021); WBEA reports analyser replacements at AMS~5 and AMS~23 in September 2024. AMS~5 has a meteorological tower with sensors at 20--90~m; its 90 minus 20~m temperature difference is the stability measure.
\end{minipage}
\end{table*}
